\section{Compiling C Control Flow and Procedures}

\subsection{Conditionals: Branching vs. Conditional Moves}

\begin{definition}
    \textbf{Compiling Control Flow:} C control flow constructs (\texttt{if-else}, loops, \texttt{switch}) are translated into assembly using \textbf{conditional jumps} (\texttt{jX}), \textbf{unconditional jumps} (\texttt{jmp}), and \textbf{conditional moves} (\texttt{cmovX}).
\end{definition}

\begin{definition}
    \textbf{Conditional Jumps vs. Conditional Moves:}
    \begin{itemize}
        \item \textbf{Conditional Jumps (\texttt{jX}):} Transfer control by modifying \texttt{\%rip} when the condition holds.
        \item \textbf{Conditional Moves (\texttt{cmovX S, D}):} Eagerly evaluate \emph{both} branches and copy \texttt{S} to \texttt{D} only if condition \texttt{X} holds, avoiding control flow jumps and branch misprediction penalties.
    \end{itemize}
\end{definition}

\begin{codeexample}
// C code: long absdiff(long x, long y) { return (x < y) ? (y - x) : (x - y); }

// 1. Branch-based:                 // 2. Conditional Move (cmov):
absdiff_branch:                     absdiff_cmov:
    cmpq  %rsi, %rdi                    movq  %rsi, %rdx
    jl    .L_then                       subq  %rdi, %rdx    // rdx = y - x
    movq  %rdi, %rax                    movq  %rdi, %rax
    subq  %rsi, %rax                    subq  %rsi, %rax    // rax = x - y
    ret                                 cmpq  %rsi, %rdi
.L_then:                                cmovl %rdx, %rax    // if x < y: rax = rdx
    movq  %rsi, %rax                    ret
    subq  %rdi, %rax
    ret
\end{codeexample}

\begin{important}
    \textbf{\texttt{cmov} cannot be used} if either branch has side effects (e.g., \texttt{x++}), is expensive to compute, or may fault when evaluated eagerly (e.g., null dereference \texttt{p ? *p : 0}).
\end{important}

\subsection{Loops and Switch Statements}

\begin{definition}
    \textbf{Loop Translations:}
    \begin{itemize}
        \item \textbf{\texttt{do-while} (Backward-Branching):} Executes the body first, then evaluates the condition at the bottom with a conditional jump backward (\texttt{jX .L\_body}).
        \item \textbf{\texttt{while} (Jump-to-Middle):} Jumps unconditionally (\texttt{jmp .L\_test}) to the test at the bottom; if true, branches back to the top of the body.
        \item \textbf{\texttt{while} (Guarded-Do):} Tests the condition initially to skip the loop if false, then runs as a backward-branching \texttt{do-while} loop.
        \item \textbf{\texttt{for (init; test; update)}:} Compiled into the same structure as a \texttt{while} loop with \texttt{init} before the loop and \texttt{update} at the end of the body.
    \end{itemize}
\end{definition}

\begin{codeexample}
// Jump-to-Middle while loop:       // Guarded-Do while loop:
    jmp   .L_test                       cmpq  $1, %rdi
.L_body:                                jle   .L_done        // Initial guard
    imulq %rdi, %rax                .L_body:
    subq  $1, %rdi                      imulq %rdi, %rax
.L_test:                                subq  $1, %rdi
    cmpq  $1, %rdi                      cmpq  $1, %rdi
    jg    .L_body                       jg    .L_body        // Backward branch
                                    .L_done:
\end{codeexample}

\begin{definition}
    \textbf{Switch Statements:}
    \begin{itemize}
        \item \textbf{Compact / Contiguous Cases $\to$ Jump Table ($O(1)$):} Stores code label addresses in an array in memory (\texttt{.rodata}). Normalizes the index ($n - \min$), checks bounds with a single unsigned comparison (\texttt{ja .L\_default}), and performs an indirect jump: \texttt{jmp *.L\_table(,\%rdi,8)}.
        \item \textbf{Sparse / Non-Contiguous Cases $\to$ Decision Tree ($O(\log N)$):} When cases are far apart, a jump table wastes memory; the compiler instead emits a binary search tree of \texttt{cmp} and conditional branches.
    \end{itemize}
\end{definition}

\subsection{Procedures and the Calling Convention}

\begin{definition}
    \textbf{The Stack and Control Transfer:} The x86-64 stack grows \textbf{downward} toward lower memory addresses, with \texttt{\%rsp} pointing to the current top of the stack.
    \begin{itemize}
        \item \textbf{\texttt{pushq S}:} $\texttt{\%rsp} \leftarrow \texttt{\%rsp} - 8$, then $\text{Mem}[\texttt{\%rsp}] \leftarrow S$.
        \item \textbf{\texttt{popq D}:} $D \leftarrow \text{Mem}[\texttt{\%rsp}]$, then $\texttt{\%rsp} \leftarrow \texttt{\%rsp} + 8$.
        \item \textbf{\texttt{call Target}:} Pushes the return address onto the stack and jumps to \texttt{Target}.
        \item \textbf{\texttt{ret}:} Pops the return address off the stack into \texttt{\%rip} to restore instruction flow.
    \end{itemize}
\end{definition}

\begin{theorem}
    \textbf{x86-64 Linux Calling Convention (System V ABI):}
    \begin{itemize}
        \item \textbf{Arguments 1--6:} Passed in \texttt{\%rdi}, \texttt{\%rsi}, \texttt{\%rdx}, \texttt{\%rcx}, \texttt{\%r8}, \texttt{\%r9} (floats in \texttt{\%xmm0}--\texttt{\%xmm7}).
        \item \textbf{Arguments 7+:} Pushed onto the stack in reverse order (7th arg at \texttt{(\%rsp)} before \texttt{call}).
        \item \textbf{Return Value:} Passed back in \texttt{\%rax}.
        \item \textbf{Caller-Saved Registers (\texttt{\%rax}, \texttt{\%rdi}, \texttt{\%rsi}, \texttt{\%rdx}, \texttt{\%rcx}, \texttt{\%r8}--\texttt{\%r11}):} Temporary registers that can be freely clobbered by the callee; the caller must save them before \texttt{call} if needed afterward.
        \item \textbf{Callee-Saved Registers (\texttt{\%rbx}, \texttt{\%rbp}, \texttt{\%r12}--\texttt{\%r15}, \texttt{\%rsp}):} Must be preserved across calls; if the callee uses them, it must push them on entry and pop them before \texttt{ret}.
    \end{itemize}
\end{theorem}

\begin{codeexample}
compute:
    pushq %rbx              // Save callee-saved %rbx
    movq  %rdi, %rbx        // Keep x safe across function call
    call  helper            // helper(x) -> result in %rax
    addq  %rbx, %rax        // return helper(x) + x
    popq  %rbx              // Restore %rbx
    ret
\end{codeexample}

\subsection{Stack Frame Optimizations and Variadic Functions}

\begin{remark}
    \textbf{Frame Pointer Omission (\texttt{-fomit-frame-pointer}):} Instead of maintaining \texttt{\%rbp} as a dedicated base frame pointer, modern x86-64 compilers compute stack frame sizes at compile time and access all local stack variables directly relative to \texttt{\%rsp}, freeing \texttt{\%rbp} for general-purpose use.
\end{remark}

\begin{definition}
    \textbf{The Red Zone:} A 128-byte region immediately below \texttt{\%rsp} (\texttt{-128(\%rsp)} to \texttt{-1(\%rsp)}) that the OS guarantees will not be overwritten by signals or interrupts. Leaf functions and small procedures can store temporary local data here without incrementing/decrementing \texttt{\%rsp}.
\end{definition}

\begin{definition}
    \textbf{Tail Call Optimization (TCO):} When a function's final action is returning the result of another function call (\texttt{return f(x);}), the compiler reuses the current stack frame and replaces \texttt{call f; ret} with a direct unconditional jump (\texttt{jmp f}), avoiding extra stack frame allocation and call/return overhead.
\end{definition}

\begin{definition}
    \textbf{Variadic Functions (\texttt{<stdarg.h>}):} Functions taking a variable number of arguments (e.g., \texttt{printf}) use compiler intrinsics to iterate through arguments passed in registers and on the stack:
    \begin{itemize}
        \item \textbf{\texttt{va\_list}:} Cursor structure tracking the current argument position.
        \item \textbf{\texttt{va\_start(ap, last\_named)}:} Initializes \texttt{ap} to point to the first variable argument.
        \item \textbf{\texttt{va\_arg(ap, type)}:} Fetches the next argument of type \texttt{type} and advances \texttt{ap}.
        \item \textbf{\texttt{va\_end(ap)}:} Cleans up the \texttt{va\_list} before the function returns.
    \end{itemize}
\end{definition}
